\section{Sistemos analizė}

Šiame skyriuje pateikiamas techninis pasiūlymas ir sistemos techninių galimybių analizė.

\subsection{Techninis pasiūlymas}

\subsubsection{Sistemos apibrėžimas}

\lipsum[1]

\subsubsection{Pagrindinis sistemos tikslas}

\lipsum[2]

\subsubsection{Sistemos pagrįstumas}

Teiginiai iš literatūros visada pagrindžiami nuoroda į šaltinį \citep{vysniauskas2006}. Kelios nuorodos sujungiamos automatiškai \citep{masiulis2007, birziska1929, gradauskas2000}, elektroniniai šaltiniai cituojami taip pat \citep{litnet2012, valiulyte2000, ieee2009}.

\lipsum[3]

\subsubsection{Egzistuojančių sprendimų analizė}

Lentelės pavadinimas rašomas virš lentelės, o tekste į lentelę būtina paminėti nuorodą prieš pačią lentelę (žr. \ref{tab:palyginimas} lentelę). Numeris ir žodis „lentelė.“ pridedami automatiškai.

\begin{table}[!ht]
\caption{Egzistuojančių sprendimų palyginimas}
\label{tab:palyginimas}
\centering
\begin{tabular}{|p{4cm}|p{3.5cm}|p{3.5cm}|p{3.5cm}|} \toprule
\textbf{Kriterijus} & \textbf{Sprendimas A} & \textbf{Sprendimas B} & \textbf{Kuriama sistema} \\ \midrule
Kriterijus 1 & Taip & Ne & Taip \\ \hline
Kriterijus 2 & Ne & Iš dalies & Taip \\ \hline
Kriterijus 3 & Taip & Taip & Taip \\ \bottomrule
\end{tabular}
\end{table}

\subsubsection{Prototipai ir pagalbinė informacija}

Paveikslo pavadinimas rašomas po paveikslu. Nuoroda į paveikslą pateikiama tekste prieš jį (žr. \ref{fig:pavyzdys} pav.).

\begin{figure}[!ht]
\centering
\includegraphics[width=3cm]{images/ktu-ikona.pdf}
\caption{Pavyzdinis paveikslas}
\label{fig:pavyzdys}
\end{figure}

\subsection{Sistemos techninių galimybių analizė}

\subsubsection{Technologijų apžvalga}

Formulės numeruojamos pagal skyrių:
\begin{equation}
  \label{eq:pavyzdys}
  P = \frac{n_{\text{padengta}}}{n_{\text{iš viso}}} \cdot 100~\%
\end{equation}
Formulėje (\ref{eq:pavyzdys}) $P$ -- padengimo procentas.

\subsubsection{Pasirinktas techninis sprendimas}

\lipsum[4]
